\section{Resonance at a regular singular point}
\label{sec:resonance}

The hypergeometric packet constructs a nonresonant connection formula. Its
remark on resonant connection formulas leaves a uniform treatment of integral
exponent differences open. Two gamma terms can each diverge in a parameter
limit, even when their sum has a well-behaved logarithmic expansion. We
construct the two local solutions directly. The classical comparison is
NIST DLMF, \S15.10(i), equations (15.10.8)--(15.10.10) and part (d);
the construction below supplies uniform numerical bounds for the
iPRA interpretation.\footnote{\url{https://dlmf.nist.gov/15.10.i}}
The same finite operation works
whenever a known solution is a monomial times a power-series unit and the
Wronskian has that form as well.

Write $K=\mathbb Q(i)$ and $|x+iy|_1=|x|+|y|$. We use $|\cdot|_1$ for
all the estimates below. A coefficient array is computed only through its
requested degree. Products of arrays use truncated convolution. For
$J\ge0$ and $R>0$ put
\[
  \|v\|_{J,R}=\sum_{j=0}^J |v_j|_1R^j .
\]
Truncated convolution has norm at most the product of the two norms: expand
the two nonnegative sums and retain the pairs with index sum at most $J$.

\begin{lemma}[Primitive of a Laurent series with one finite principal part]
\label{lem:res-primitive}
Suppose every prefix of $g$ has $\|g\|_{J,R}\le M$. For $k\in\mathbb N$,
the expression
\[
 H(t)=g_k\log t+
 \sum_{\substack{j\ge0\\j\ne k}}
       \frac{g_j}{j-k}t^{j-k}
\]
gives a primitive of $t^{-k-1}g(t)$ on a named logarithm chart with
$0<|t|_1\le R/2$. Keeping $j\le k+N$ gives error at most
$MR^{-k}2^{-N-1}$. The logarithm is part of the exact expression even at
cutoff $N=0$.
\end{lemma}
\begin{proof}
Every omitted exponent is positive. For every subsequent finite tail,
multiply and divide its terms by $R^j$, use $1/(j-k)\le1$, and bound
$(|t|_1/R)^{j-k}$ by $2^{-N-1}$. The resulting bound is independent of
the second cutoff. Formal differentiation sends the exceptional summand to
$g_k/t$, and sends every other summand to $g_jt^{j-k-1}$. At a requested
point the negative powers are finite and require a supplied apartness
bound. Derivatives of the positive tail converge with explicit clocks too:
the $d$th derivative, $d=1,2$, costs a factor at most
$R^{-d}2^d(j-k)^d$ against the bound $R^{-k}2^{-(j-k)}$ for
$j-k\ge d$. The elementary bound $n^2\le 64\,2^{\lfloor n/2\rfloor}$
for $n\ge1$ follows by checking $1\le n\le8$ and stepping by two.
Thus twice the requested bit precision plus an integer offset for the
displayed constants suffices. Termwise differentiation here can equivalently
be read as a finite Taylor estimate: for polynomials expand
$(t+h)^n-t^n-n t^{n-1}h$ by the binomial identity and bound the remainder
by $n(n-1)|h|_1^2(|t|_1+|h|_1)^{n-2}/2$. On a slightly larger disk
(e.g. $3R/4$) the weighted derivative tails still decay geometrically
(group four powers to use $(3/4)^4<1/2$). The rational inequalities pass
to names at any prescribed precision.
\end{proof}

\subsection{A basis uniform in the signed exponent difference}

Fix Gaussian rationals $a,b,c$ and a signed integer $m=c-a-b$. Use
$t=1-z$, $k=|m|$, $s=\min(0,m)$ and $r=\max(0,m)$. The following formulas
also work when upper parameters terminate. Let $h_0=B_0=u_0=1$, and set
\[
\begin{aligned}
 h_{j+1}&=h_j\frac{(a+r+j)(b+r+j)}{(j+1)(j+1+k)},&
 B_{j+1}&=B_j\frac{c+j}{j+1},\\
 u_n&=-\sum_{j=1}^n h_j u_{n-j}\quad(n>0),&
 g&=B*u*u .
\end{aligned}
\]
Here $B$ is the series for $(1-t)^{-c}$ at zero. This notation can be
defined by its displayed recurrence, equivalently by $(1-t)B'=cB$ and
$B(0)=1$. Use the rational radius
\[
 R=\big[4\max\{1,(1+|a+r|_1)(1+|b+r|_1),1+|c|_1\}\big]^{-1}.
\]

\begin{theorem}[Local resonant solutions and logarithmic coefficient]
\label{thm:res-gauss}
On $0<|t|_1\le R/2$ with a specified logarithm branch, a fundamental pair
for Gauss's equation in the coordinate $t$ is
\[
 U=t^r h(t),\qquad V=U H(t),
 \qquad U V'-U' V=t^{m-1} B(t),
\]
where $H$ is given by Lemma~\ref{lem:res-primitive}. All coefficient and
tail algorithms are uniform primitive recursive functions of the rational
parameters, the signed resonance and the requested precision, with the
point and branch guards described below. With $(x)_k=\prod_{j<k}(x+j)$
and $(x)_0=1$, the logarithmic coefficient in $H$ is
\[
 \boxed{\displaystyle
 g_k=\frac{(-1)^k(a+s)_k(b+s)_k}{(k!)^2}.}
\]
Consequently a change of logarithm by $2\pi i$ sends $V$ to
$V+2\pi i g_k U$. For rational parameters this also decides exactly when
the logarithmic term disappears.
\end{theorem}

\begin{proof}
A substitution in the equation gives
\[
 L y=t(1-t)y''+[1-m-(a+b+1)t]y'-aby=0 .
\]
Inserting $t^e v(t)$ and dividing by $t^e$ contributes a term
$e(e-m)v/t$. For $e=r$ or $s$ this vanishes. The resulting Gauss operator
has upper parameters $a+e,b+e$ and derivative constant $1-m+2e$.
For $e=r$ this constant is $1+k$. Coefficient comparison gives exactly
the recurrence for $h$ above, hence $LU=0$.

For every prefix the coefficient ratios bound
$|h_j|_1 R^j, |B_j|_1 R^j\le 4^{-j}$. For example
$|a+r+j|_1\le(1+|a+r|_1)(j+1)$, and the denominator in the recurrence
for $h$ is at least $(j+1)^2$. Therefore $\|h-1\|_{J,R}\le1/3$.
In the algebra of polynomials modulo $t^{J+1}$,
$u=\sum_{d=0}^{J}(1-h)^d$, since the error term has degree at least
$J+1$. This proves the displayed inverse recurrence,
$\|u\|_{J,R}\le3/2$ and $\|g\|_{J,R}\le(4/3)(3/2)^2=3$.
Replacing $c$ by $-c$ gives an array $B_{-c}$ with the same bound
$4/3$. The formal derivative of $B B_{-c}$ is zero and its constant is
one, by the first-order equations and cancellation of $1-t$ in truncated
coefficient calculations. Thus $B^{-1}=B_{-c}$.

In the normalized ODE the coefficient of $y'$ is
$p=(1-m)/t-c/(1-t)$. Consequently $W=t^{m-1}B$ satisfies $W'=-pW$.
From $r-s=k$ and $r+s=m$ we get $W/U^2=t^{-k-1}g$. Defining
$V=UH$ gives $UV'-U'V=U^2 H'=W$. Differentiating this equation gives
$U(V''+pV'+qV)=0$ because $U''+pU'+qU=0$, where
$q=-ab/[t(1-t)]$. Multiplication by $t^{-r}u$ cancels $U$.
The inverse determinant has the explicit name $t^{1-m}B_{-c}$. This
constructs a fundamental matrix with an inverse, even if $g_k=0$.
For any solution germ presented by a convergent Taylor expansion at an
ordinary point, multiplying its initial jet by this inverse gives its
unique two coordinates. The jet uniqueness and transport estimates below
prove equality of the resulting solutions throughout any specified chain
of overlapping charts.

It remains to identify the logarithm. If $k=0$, $g_0=1$. For $k>0$, write
$t^{-s}V=Q+g_k t^k h\log t$. Here $Q$ is an ordinary power series, since
$t^k$ cancels the finite principal part, and $Q_0=-1/k$.
Put $A=a+s$, $D=b+s$. The transformed equation has derivative constant
$1-k$. Applying its operator to $t^k h\log t$, the logarithmic part
vanishes. The remainder is
$2(1-t)(t^k h)' +[-k/t-(A+D)]t^k h$, whose first term of degree $k-1$
has coefficient $k$. For $1\le j<k$ coefficient comparison at $j-1$ now
gives
\[
 j(j-k)Q_j=(A+j-1)(D+j-1)Q_{j-1}.
\]
At $j=k$ the coefficient of $Q_k$ is zero and
$kg_k=(A+k-1)(D+k-1)Q_{k-1}$. Multiply the first $k-1$ recurrences,
using $\prod_{j=1}^{k-1}j(j-k)=(-1)^{k-1}((k-1)!)^2$.
Iteration uses division only by the nonzero integers $j(j-k)$, including
when an upper rising factorial vanishes. This proves the boxed identity.
Changing only the logarithm in $H$ gives the asserted shear.
\end{proof}

\paragraph{Precision data.}
The finite negative-power sum uses rational arithmetic at Gaussian-rational
points. At named points, a lower bound $|t|_1\ge\delta>0$ gives
$|1/t|_1\le 2/\delta$, since $(|x|+|y|)^2\le 2(x^2+y^2)$.
For a branch supply an evaluation procedure for $\log t$ with its error
clock and a bound on its norm. A chart based at $x$ translates by
$\log(x+u)=\log x+\sum_{j\ge1}(-1)^{j+1}(u/x)^j/j$ for $|u/x|_1<1$,
with the geometric tail from the supplied radius margin. This fixes the
derivative and the overlap convention. Positive rational $t$ has a canonical real logarithm from
$\log t=2\sum_{j\ge0} v^{2j+1}/(2j+1)$, $v=(t-1)/(t+1)$, with tail
bounded by $2|v|^{2J+1}/(1-|v|^2)$ after $j<J$.
For the power-series tails take $M=3$ in Lemma~\ref{lem:res-primitive}.
A cutoff for accuracy $2^{-p}$ is, for example,
$N=p+\lceil 3 R^{-k}\rceil+2$ for $H$ before multiplication by $U$.
All intermediate size bounds can be propagated through finite sums and
products: if $|x|_1\le A, |y|_1\le D$ and approximants have errors
at most $\epsilon\le1$, their product error is at most
$(A+D+1)\epsilon$. This gives rational instruction-by-instruction budgets
for $U,V,U',V'$ and for the inversions above. Every induction may fix the
requested precision and a final degree $J$ first. The needed prefix
recurrences, truncated product identities, and prefix norm inequalities
are then Boolean folds on rational arrays.

\subsection{Matching the solution at zero without gamma limits}
Assume $c\notin\{0,-1,-2,\ldots\}$, a decidable condition on $K$. The
ordinary series $F(z)={}_2F_1(a,b;c;z)$ then has coefficients
$f_0=1$ and $f_{n+1}(n+c)(n+1)=f_n(n+a)(n+b)$.
Set $t_0=R/4$, $z_0=1-t_0$. Use the real logarithm at $t_0$.
With $Y=F(z_0)$, $D_0=-F'(z_0)$ and $W=UV'-U'V$, define
\[
 \alpha=(YV'-D_0 V)/W,\qquad \beta=(U D_0-U'Y)/W ,
\]
where the right sides are evaluated at $t_0$.

\begin{proposition}[Connection coefficients as rapid names]
On the overlap containing the positive interval $0<t\le R/2$,
the continuation of $F(1-t)$ from $t_0$ is $\alpha U+\beta V$, using
the logarithm continued from its real value at $t_0$ along the same
charts. The constants are uniformly PR-computable names from
$a,b,c,m$ alone.
\end{proposition}
\begin{proof}
Here is a clock even when $z_0$ is near one. Put
$E=\max(1,|a|_1,|b|_1,|c|_1)$,
$n_0=\lceil32 E/t_0\rceil$, and $q=1-t_0/2$. For $n\ge n_0$ put
$x=E/n\le 1/2$. The ratio multipliers of $f_n z_0^n$ and
$n f_n z_0^{n-1}$ both have norm at most
$z_0(1+x)^2/(1-x)$. For the derivative the factor $(n+1)/n$ cancels the
coefficient denominator $(n+1)/n$; the series itself has the smaller
bound. The reciprocal estimate follows by multiplying
$(1+c/n)^{-1}=1-(c/n)(1+c/n)^{-1}$ and taking norms. For $x\le1/2$,
$(1+x)^2/(1-x)\le1+7x$. Thus both norms are $\le q$. This argument uses
multipliers, so a vanishing term creates no division by zero. Let $M_0$ be
one plus the sum of the two term norms at $n_0$. All subsequent tails
after $N$ further terms are bounded by $M_0q^N/(1-q)$. For an explicit
loose cutoff, $q^N\le [1+N(1-q)]^{-1}$ by the finite binomial inequality.
Hence $N=\lceil 2^p M_0/(1-q)^2\rceil+1$ suffices at precision $p$.
The same inequality supplies a clock for the logarithm at $t_0$.
The previous theorem supplies the other four entries and $W^{-1}$.
Matrix multiplication gives equality of values and first derivatives
at $t_0$ exactly as names.

For completeness, agreement of jets at an ordinary point $x$ determines
the expansion there. After translation, the coefficient of the still
unknown Taylor coefficient of degree $j+2$ in the differential equation
is $(j+1)(j+2)x(1-x)$; all other terms involve earlier coefficients.
Hence the two coefficient arrays agree to every prescribed degree by
finite recurrence. Reexpansions used here exist by the binomial formula:
for a series having bounded weighted prefixes at radius $S$, expanding
$(x+u)^n$ is absolutely majorized by $(|x|_1+|u|_1)^n$. A strict gap
to $S$ gives the geometric truncation clock; locally the logarithm
translates by the Taylor series for $\log(1+u/x)$ for $|u/x|_1<1$
with its geometric tail. For the original Gauss series at a positive
point use the ratio bound above at any larger radius below one.
The Laurent powers translate by the binomial series or finite
multiplication of the local inverse. Thus on the real segment between
$t_0$ and a specified $t\in(0,R/2]$, both expansions are valid on disks
of radius $\min(t,t_0,R/4)/8$ about every segment point (for a named
$t$ use its supplied lower bound). Here are finite
bounds for the dependence on these jets. Put $d=\min(t,t_0,R/4)$
(or a positive rational lower bound), $S=d/4$, and write the differential
equation as the first-order system for $(y,y')$ with matrix
\[
 P(t)=
 \begin{pmatrix}
 0&1\\
 ab[1/t+1/(1-t)]&(m-1)/t+c/(1-t)
 \end{pmatrix}.
\]
At a segment center $x$, summing the geometric expansions of $1/t$ and
$1/(1-t)$ out to radius $S$ bounds their prefix norms by
$1/(x-S)\le 2/d$ and $1/(1-x-S)\le2$, respectively. Thus
$A=1+|ab|_1(2/d+2)+2|m-1|/d+2|c|_1$ bounds the sum of all entry norms.
Set $\rho=\min(S,1/(2A),d/8)$. For initial vector $b_0$, the recurrence
$(j+1)b_{j+1}=\sum_{i=0}^j P_i b_{j-i}$ gives
\[
 \|b\|_{J,\rho}\le |b_0|_1+\rho A\|b\|_{J-1,\rho}\le 2|b_0|_1 ,
\]
by finite induction (using the sum of coordinate norms). A transfer of
length $\le\rho/2$ has operator norm at most two, and truncating at $J$
costs at most $2|b_0|_1 2^{-J-1}$. Divide the segment into
$L=1+\lceil2|t-t_0|/\rho\rceil$ such steps. Initial errors grow by
at most $2^L$. If each transfer adds error $\epsilon$, the aggregate
added error is at most $(2^L-1)\epsilon$. Thus set
$\epsilon=2^{-p-L-2}$, and bound intermediate exact jets by
$2^L|b_0|_1$ when choosing cutoffs. These rational bounds give an
induction on errors at fixed precision in place of induction on equality
of entire analytic functions. The same jet argument extends
to specified overlap charts by finite chains with positive margins.
\end{proof}

\paragraph{Worked case.}
For $(a,b,c)=(1,1,3)$ the resonance is $m=1$,
$h=(1-t)^{-2}$ and $g=1-t$. Our basis is
$U=t/(1-t)^2$ and $V=U(-1/t-\log t)$.
The coefficients of the series at zero are $2/((n+1)(n+2))$,
so summing the partial fractions using the logarithm series gives
\[
 F(1-t)=\frac{2[(1-t)+t\log t]}{(1-t)^2}=-2U-2V .
\]
The sign of the shear and the normalization of the principal part can
both be checked on this example. For a resonance with trivial logarithm,
$(a,b,c)=(0,1,3)$ gives $k=2$ and $(a)_2=0$.

\paragraph{Formalization targets and scope.}
The repository's \texttt{HypergeometricRatio} proves cross-multiplied
identities with natural parameters. \texttt{HypergeometricOdeFormal}
already derives the Wronskian derivative in its formal differential
algebra. A Lean implementation can encode each of $h,B,u,g$ as a
degree-$J$ array of pairs of signed rationals and use the recurrences
above as bounded loops. First derive multiplication and inversion modulo
$t^{J+1}$, prefix norm submultiplicativity, and the $4^{-j}$ bounds in
the object calculus. Next encode a finite Laurent jet by a signed lower
exponent, a degree bound, and its regular and logarithmic arrays, with
$D(t^n\log t)=nt^{n-1}\log t+t^{n-1}$. To verify a coefficient of
$Ly$ through exponent $n$ request the coefficients of $y$ through $n+2$,
also allowing for the displayed monomial shifts. The boxed formula is
a rational-array recurrence with finite product invariants; the Wronskian
and basis identities reduce coefficientwise to convolution. The final
stage uses the clocks above to translate each requested identity into
rational inequalities. The included \texttt{Resonance.lean} checks six host-field
lemmas in Lean 4.32.0, including the coefficient recurrence telescope,
the exceptional-index identity and Wronskian matching. Its field
identities are reusable algebraic inputs. Encoding arrays and deriving
their norm and analytic bounds in object iPRA remain formalization work.
The exact check in \texttt{checks/resonance.py} tests 33 parameter triples
with $-5\le m\le5$ through degree 13, including terminating and
nonreal parameters. The companion coefficients and monodromy shear
are local. The matching formula gives numerical names for the two
connection constants. An expression for these constants in gamma and
digamma values or a classification of global monodromy requires additional
theorems.

The accompanying \texttt{ResonanceObject.lean} carries the multiplicative
telescope into ordinary iPRA. For fixed PR term bodies $a_j(x),b_j(x),v_j(x)$,
its hypotheses are the object formulas $b_jv_{j+1}=a_jv_j$ for all $j<N$
and $\ell g=a_Nv_N$. Its closed conclusion gives
\[
 \left(\prod_{j<N}b_j\right)\ell g=\left(\prod_{j\le N}a_j\right)v_0 .
\]
The horizon $N$ is an object variable. This lemma uses natural-number
products. Using it for general Gaussian coefficients requires signed
arithmetic and denominator bookkeeping. A second closed derivation handles
the first positive resonance: from the cleared recurrence for $g_1=p-q$
with natural $a,b$, it proves $p+ab=q$.
